% Approximate-composition lemma (original Proposition 6), clean statement, proof and tightness.
% Terminal 2 (claims role, new numbering). Standalone: pdflatex COMPOSITION_LEMMA.tex.
% For the manuscript, copy the lemma, proof and remark into the appendix; the body sentence already
% present at main.tex@8df7527c:171-175 is consistent with part (i).
\documentclass{article}
\usepackage{amsmath,amssymb,amsthm}
\newtheorem{lemma}{Lemma}
\newtheorem*{remark}{Remark}
\DeclareMathOperator{\tr}{tr}
\begin{document}

\paragraph{Setting.} Let $h_1,\dots,h_k$ be random vectors ($h_p\in\mathbb R^{d_p}$) and $A\in\mathbb R^{c}$ a
target (for a categorical attribute, the centred one-hot vector), all with finite second moments, on one
distribution. That distribution may be the empirical distribution of an evaluation sample, in which case
every covariance below is the sample covariance. Write $\Sigma_p=\operatorname{Cov}(h_p)$,
$\Sigma_{pq}=\operatorname{Cov}(h_p,h_q)$, $\Sigma_{pA}=\operatorname{Cov}(h_p,A)$ and
$\Sigma_A=\operatorname{Cov}(A)$, with $\tr\Sigma_A>0$. For a view $h$ with $\Sigma_h\succ0$ and
$\tau\ge 0$ define the affine least-squares score
\[
R^2_\tau(h;A)=\frac{\tr\big(\Sigma_{hA}^{\top}(\Sigma_h+\tau I)^{-1}\Sigma_{hA}\big)}{\tr\Sigma_A},\qquad R^2=R^2_0 .
\]
Assume $\Sigma_p\succ0$ for every $p$. Let $D=\operatorname{blkdiag}(\Sigma_1^{1/2},\dots,\Sigma_k^{1/2})$ and
let $R=D^{-1}\Sigma_H D^{-1}$ be the block-whitened correlation of $H=(h_1;\dots;h_k)$, so $R_{pp}=I$ and
$R_{pq}=\Sigma_p^{-1/2}\Sigma_{pq}\Sigma_q^{-1/2}$.

\begin{lemma}[Approximate composition]
\begin{enumerate}
\item[(i)] If $R\succ0$, then $R^2(H;A)\le \sum_p R^2(h_p;A)\,/\,\lambda_{\min}(R)$.
\item[(ii)] If $R$ is singular, then (i) holds for the pseudo-inverse score
$\tr(\Sigma_{HA}^{\top}\Sigma_H^{+}\Sigma_{HA})/\tr\Sigma_A$, with $\lambda_{\min}(R)$ replaced by the smallest
positive eigenvalue $\lambda_+(R)$.
\item[(iii)] For every $\tau\ge0$ and $R\succ0$: $R^2_\tau(H;A)\le \sum_p R^2_\tau(h_p;A)\,/\,\lambda_{\min}(R)$,
with the \emph{same} $\tau$ in every score.
\item[(iv)] Parts (i)--(iii) hold for each target column separately (the per-class score), and for every
contrast $a^\top A$. Hence they hold for the largest squared canonical correlation between $H$ and $A$.
\end{enumerate}
\end{lemma}

\begin{proof}
Let $M_p=\Sigma_p^{-1/2}\Sigma_{pA}$ and $M=(M_1;\dots;M_k)=D^{-1}\Sigma_{HA}$. Then
$R^2(h_p;A)=\|M_p\|_F^2/\tr\Sigma_A$. Since $\Sigma_H=DRD$,
$R^2(H;A)=\tr(M^\top R^{-1}M)/\tr\Sigma_A\le \|M\|_F^2/(\lambda_{\min}(R)\tr\Sigma_A)$, because
$R^{-1}\preceq\lambda_{\min}(R)^{-1}I$. Finally $\|M\|_F^2=\sum_p\|M_p\|_F^2$. This proves (i).

For (ii): if $Ru=0$ then $\operatorname{Var}(u^\top D^{-1}H)=0$, so $u^\top M=0$. The columns of $M$ therefore
lie in $\operatorname{range}(R)$, where $R^{+}\preceq\lambda_+(R)^{-1}I$.

For (iii): write $\Sigma_H+\tau I=D(R+\tau D^{-2})D$ and $\Sigma_p+\tau I=\Sigma_p^{1/2}(I+\tau\Sigma_p^{-1})\Sigma_p^{1/2}$.
The per-view score is then $\tr(M_p^\top(I+\tau\Sigma_p^{-1})^{-1}M_p)/\tr\Sigma_A$. Put
$K_d=I+\tau D^{-2}$. Because the diagonal blocks of $R$ are identities, $\tr R=\dim R$, so
$\lambda:=\lambda_{\min}(R)\le1$. Hence $R+\tau D^{-2}\succeq\lambda I+\tau D^{-2}\succeq\lambda K_d$, which gives
$(R+\tau D^{-2})^{-1}\preceq\lambda^{-1}K_d^{-1}$. Since $K_d$ is block diagonal,
$\tr(M^\top K_d^{-1}M)=\sum_p\tr(M_p^\top(I+\tau\Sigma_p^{-1})^{-1}M_p)$.

For (iv): apply (i)--(iii) with $A$ replaced by a single column, or by $a^\top A$, and take the supremum
over unit-variance contrasts: each summand is at most that view's largest squared canonical correlation.
\end{proof}

\paragraph{Tightness.} Let $V,S$ be independent, centred, unit-variance, and set $h_1=V+\delta S$,
$h_2=V-\delta S$ with $\delta\ne0$. Then $R^2(h_p;S)=\delta^2/(1+\delta^2)$, $R_{12}=(1-\delta^2)/(1+\delta^2)$,
$\lambda_{\min}(R)=2\delta^2/(1+\delta^2)$, and the bound equals $1=R^2(H;S)$, because $h_1-h_2=2\delta S$. The
ill-conditioning is explicit: $\operatorname{cond}(\Sigma_H)=\delta^{-2}$.

\begin{remark}[Mapping to the implemented statistic and missing assumptions]
\begin{itemize}
\item \emph{The repository statistic.} It is an in-sample ridge score with gram
$H_c^\top H_c+\lambda I$, i.e.\ $\tau=\lambda/n$ in covariance units, on the evaluation sample. Part (iii)
applies to it exactly, with the \emph{sample} covariances and the same $\lambda$ for every view and for
the concatenation. It is a statement about that sample, not about the population or a held-out attacker.
\item \emph{Collapsed representations.} Some views have singular or near-singular $\Sigma_p$ (effective
rank $\approx3$ in the original encoder line). Then $R$ is undefined or $\lambda_{\min}(R)$ is tiny, and
the bound is vacuous unless (ii) is used on the supported subspace. An application must report
$\lambda_{\min}(R)$ or $\lambda_+(R)$ and the eigenvalue threshold used.
\item \emph{Max-single-class versus all-class-contrast.} The pooled score is a prevalence-weighted
average of the per-class scores. The implemented dominant-axis score is the maximum per-class score,
which is itself at most the top squared canonical correlation. The lemma bounds each of these
separately (part (iv)). A bound on the per-class maxima does not bound a class contrast.
\item \emph{Scope.} The lemma is linear and second-order. It implies nothing about nonlinear recovery:
two views each with zero covariance with $S$ can determine $S$ exactly.
\item \emph{Status.} This is an elementary Rayleigh--Ritz bound, not a new theorem of the method. Part
(iii) with a regularised premise corrects an earlier over-cautious note in the claims audit (T3-12),
which said the premise had to be unregularised.
\end{itemize}
\end{remark}
\end{document}
